

% \subsection{Poisson-Neumann Data}

% \begin{equation}
% \partial_{\bm{\nu}} \bm{\psi}=\nu \mathbf{I}(a \mathbf{1}+b \mathbf{J}) \bm{\nu}
% \end{equation}
% \begin{equation}
% \partial_{\bm{\nu}} \bm{\psi}=\nu \mathbf{I}\left(a \mathbf{1}+b\left[\mathbf{i}^{\times}\right]\right) \bm{\nu}
% \end{equation}

% with the two scalar quadratics $a, b$ extracted from $R$ as the components of the deviatoric part \(\operatorname{dev} \bm{R}\).

% Let $D:=\operatorname{dev} \mathbf{I}=\mathbf{I}-\frac{1}{2}(\mathbf{I}: \mathbf{1}) \mathbf{1}$ (constant over the section).
% Form the two symmetric deviators without naming principal directions:

% $$
% \mathbf{U}:=\frac{D}{2 D:D}, \quad \mathbf{~B}:=\frac{\operatorname{sym}\left(\left[\mathbf{i}^{\times}\right] D\right)}{D: D}=\frac{\left[\mathbf{i}^{\times}\right] D-D\left[\mathbf{i}^{\times}\right]}{D: D} .
% $$
% Define:
% $$
% a=\frac{1}{2} \operatorname{dev} \bm{R}: \mathbf{U}_I, \quad b=\operatorname{dev} \bm{R}: \mathbf{V}_I .
% $$

% Then set
% $$
% a=\operatorname{dev} \bm{R}: \mathbf{U}, 
% \quad 
% b=\operatorname{dev} \bm{R}: \mathbf{V} .
% $$
% With these definitions,

% $$
% \mathbf{\Gamma}=\nu \operatorname{dev} \bm{R}=\nu\left(a \mathbf{1}+b\left[\mathbf{i}^{\times}\right]\right) = -\frac{\nu}{2}(\bm{Q} - \bm{R})
% $$


\subsection{Saint-Venant's Problem}

Our primary objective is to obtain a model that reproduces the force-displacement response of Saint-Venant's tip loaded cantilever. 
To this end, it is useful to first consider the case of torsion warping. In this case one warping mode is introduced with amplitude proportional to the rate of twist $\vartheta'$. 
The subtlety here is that even when warping is induced by $\vartheta'$, the conjugate resultant $\tilde{T}$ remains essentially equivalent to the classical torsion, $T=\mathbf{i}\cdot\int \bm{r}\times\bm{\tau}d\Omega$. 
To see this, consider the strains for the model (under pure torsion) are $\bm{\varepsilon} = \vartheta'(\mathbf{i}\times\bm{r} + \nabla\tilde{\varphi})$, where $\tilde\varphi$ is the SV warping function. The conjugate resultant $\tilde{T} = \int (\mathbf{i}\times\bm{r} + \nabla\tilde{\varphi}) \cdot G (\mathbf{i}\times\bm{r} + \nabla\tilde{\varphi}) d\Section$ reduces to \emph{exactly} $T = \int \mathbf{i}\times\bm{r} \cdot G (\mathbf{i}\times\bm{r} + \nabla\tilde{\varphi}) d\Section$, because by definition of the SV function $\tilde\varphi$, we remarkably have $\int \nabla\tilde\varphi \cdot (\mathbf{i}\times\bm{r} + \nabla\tilde{\varphi}) d\Section \equiv 0$. 
Recall that $\tilde\varphi$ is defined by
\begin{equation}
\begin{array}{rll}
\operatorname{div}_G\left(\left.\nabla \WarpSVP+\mathbf{i} \times\bm{r}\right.\right) &=0 & \text { in } \PlaneManifold, \\
\nabla_{\bm{\nu}} \WarpSVP &=-\mathbf{i}\times\bm{r} \cdot \bm{\nu} & \text { on } \PlaneBoundary, \\
\int_{\PlaneManifold} \WarpSVP \volE &= 0
\end{array}
\label{eq:warp-sv-bvp}
\end{equation}

\begin{equation}
\bm{p} \triangleq \frac{1}{E}\mathbf{J}^{-1}\bm{q}
\end{equation}

\begin{equation*}
\begin{cases}
    \sigma = E \, \bm{r}\cdot\mathbf{D}\bm{\gamma}\\
    \bm{\tau} = G\left(\nabla \bm{\psi}^{\mathrm{t}} + \bm{A}_{d}\right) \mathbf{D}\bm{\gamma}
\end{cases}
\end{equation*}


\begin{equation*}
    \operatorname{div} \bm{\tau} = 
    G(\Delta (\bm{\psi} \cdot\mathbf{D}\bm{\gamma}) - 2 \nu \, \bm{r}\cdot \mathbf{D}\bm{\gamma})
\end{equation*}

Let 
\begin{equation*}
\bar{\nu} = \frac{\nu}{1+\nu} = 2\nu \frac{G}{E}
\end{equation*}
and note:
\begin{equation*}
\frac{G}{E} = \frac{1}{2(1+\nu)}
\end{equation*}
% Define 
% \begin{equation}
% \mathbf{S}:=\operatorname{Off}(\bm{r}-\mathbf{Q})=2 y z\left(\mathbf{e}_y \otimes \mathbf{e}_z+\mathbf{e}_z \otimes \mathbf{e}_y\right) .
% \end{equation}

\subsection{Warping Functions}
Several distinct warping functions have been investigated in the literature. 
The can generally be categorized by:
\begin{itemize}
    \item Nature (Neumann vs Dirichlet)
    \item Forcing (Laplace vs Poisson)
    \item Boundary (shift by harmonic gauge)
    \item Splitting
    \item Scaling
\end{itemize}

\subsubsection{Splitting}
Recall that given any harmonic function $u, \nabla u$ is solenoidal since
$$
\nabla \cdot(\nabla u)=\nabla^2 u=0 .
$$
Consider a \emph{Helmholtz decomposition} of the shear stress \(\bm{\tau}\):
\begin{equation}
\boldsymbol{\tau}=\underbrace{\nabla \phi^{\text {irr }}}_{\text {curl }=0}+\underbrace{\mathbf{i}\times \nabla \phi^{\text {sol }}}_{\text {div }=0},
\label{eq:helmoltz-shear-stress}
\end{equation}

\begin{equation}
\begin{aligned}
 \tau &= \nabla \bar{\Psi}_1 + \mathbf{i}\times\nabla\Psi_2 \\
&=\underbrace{\nabla q}_{\text {irrotational }}+\underbrace{\mathbf{i}\times \nabla \Psi}_{\text {solenoidal }} \\
&=\nabla \bar{\Psi}+\mathbf{i}\times \nabla p
\end{aligned}
\end{equation}
Note there is non-uniqueness in the split of \eqref{eq:helmoltz-shear-stress}. In 2-D any shear field \(\bm{\tau}\) is invariant under the ``harmonic gauge"
$$
\phi^{\mathrm{irr}} \rightarrow \phi^{\mathrm{irr}}+h, \quad \phi^{\mathrm{sol}} \rightarrow \phi^{\mathrm{sol}}+g \quad \text { with } \quad \nabla h=-\mathbf{i} \times \nabla g,
\quad \Delta h=\Delta g=0 .
$$
\begin{equation*}
\boldsymbol{\tau}=\underbrace{\nabla \chi}_{\text {irrotational }}+\underbrace{\mathbf{i} \times \nabla \omega}_{\text {solenoidal }}+\underbrace{\boldsymbol{\eta}}_{\text {harmonic }}
\end{equation*}

\begin{equation}
\begin{array}{lccccccccc}
\mathrm{ker} \text { curl }= &  && \mathrm{HK} &\oplus& \mathrm{CG}& \oplus& \mathrm{HG} &\oplus& \mathrm{GG} \\
\mathrm{im } \text { grad }= &  &&             &      & \mathrm{CG} &\oplus& \mathrm{HG} &\oplus& \mathrm{GG} \\
\mathrm{im } \text { curl }= & \mathrm{FK} &\oplus& \mathrm{HK} &\oplus& \mathrm{CG} &\\
\mathrm{ker} \text { div  }= & \mathrm{FK} &\oplus& \mathrm{HK} &\oplus& \mathrm{CG} &\oplus& \mathrm{HG}
\end{array}
\end{equation}
\begin{equation}
\begin{array}{ll}
\mathrm{FK}=\text{fluxless knots } &=\{\nabla \cdot V=0, V \cdot n=0, \text { all interior fluxes are } 0\} \\
\mathrm{HK}=\text{harmonic knots } &=\{\nabla \cdot V=0, \nabla\times V=0, V \cdot n=0\} \\
\mathrm{CG}=\text{curly gradients} &=\{V=\nabla\varphi, \nabla\cdot V=0, \text { all boundary fluxes are } 0\} \\
\mathrm{HG}=\text{harmonic gradients}&=\{V=\nabla\varphi, \nabla\cdot V=0, \varphi \text{ is locally constant on } \partial \Section\} \\
\mathrm{GG}=\text {grounded gradients }&=\left\{V=\nabla \varphi,\left.\phi\right|_{\partial \Section}=0\right\} .
\end{array}
\end{equation}


\subsubsection{Framework}

\begin{enumerate}
    \item Obtain classical warping functions \(\bm{\phi}:\Section\rightarrow\mathbb{R}^2\), the gradient \(\nabla \bm{\phi}\), and the corresponding stress correction operator \(\bm{A}_{\psi}:\Section\rightarrow\mathbb{R}^{2\times 2}\) through a classical Neumann problem.
    \item Obtain the optimally gauged nominal warping functions \(\hat{\bm{\phi}}\) by solving:
    \begin{equation}
    \hat{\bm{\phi}}=\arg \min _\phi \int_{\Section}\left\|\nabla_{\bm{u}} { \bm \phi}-(\nabla_{\bm{u}} \bm{\phi} + \bm{A}_{\phi}\bm{u})\right\|^2 \,\mathrm{d} A \Longleftrightarrow 
    \begin{cases}
    -\Delta \hat{\bm \phi}=\nabla \cdot (\nabla \bm{\phi} + \bm{A}_{\phi}) & \text { in } \Section, \\ 
    \partial_{\bm{\nu}} \hat{\bm\phi}=\bm{\nu} \cdot (\nabla \bm{\phi} + \bm{A}_{\phi}) & \text { on } \partial \Section, \\ 
    \int_{\Section} \hat{\bm \phi}\, \mathrm{d} A=\mathbf{0} & \text {uniqueness. }
    \end{cases}
    \end{equation}
    \item Extract the projected interpolation shapes \(\bm{\varphi} = \bm{S}\Pi[\hat{\bm{\phi}}]\) \citep{simo1982bifurcation}.
\end{enumerate}


\subsection{Neumann Problems}
For problems which are only well-posed in the centroidal coordinate system, the problem is referred to \(\bar{\Section}\).

\begin{equation}
\begin{aligned}
& \Phi_{\mathrm{p}}=-\left(\chi_{\mathrm{p}}+f_{\mathrm{p}}\right)+\operatorname{poly}(M) . \\
& \bar{\Psi}_{\mathrm{p}}=G \chi_{\mathrm{p}}+\operatorname{poly}(M) . \\
& \varphi_{\mathrm{p}}=\alpha_{\mathrm{p}} \Pi_{\mathrm{p}}\left[\Phi_{\mathrm{p}}\right] .
\end{aligned}
\end{equation}



\subsubsection{Laplace-Neumann I (L0-N2, Love)}
\cite[Art. 229]{love1944treatise}:
Love’s torsionless boundary–value problem reads
\begin{equation}
\left\{
\begin{aligned}
\Delta\,\bm{\chi}&=\mathbf{0} &&\text{in }\bar{\Section},\\[2pt]
\partial_{\bm{\nu}}\boldsymbol{\chi}&= -\,\mathbf A(\bm{r})\,\bm{\nu}
&&\text{on }\partial\bar{\Section},\\[2pt]
\displaystyle\int_{\Section}\boldsymbol{\chi}\mathrm{~d}A&=\mathbf{0} &&\text{uniqueness}, \\
\int_{\Section}\bm{r}\mathrm{~d}A&=\mathbf{0} &&\text{compatibility}.
\end{aligned}
\right.
\label{eq:love-flexure-function}
\end{equation}
with
\begin{equation}
\mathbf A(\bm{r})\;=\;\tfrac12(\bm{r}\cdot\bm{r})\,\mathbf{1}
\;+\;(\nu-1)\,u_I\,\mathbf{U}_{\rm I}
\;+\;(2\nu+1)\,v_I\,\mathbf{V}_{\rm I}.
\end{equation}
The units of \(\LoveWarp\) are length$^3$. 
\eqref{eq:love-flexure-function} is only well-posed in a centroidal coordinate system. \cite{sokolnikoffMathematicalTheoryElasticity1946} uses the symbol \(\Phi\) for \(\LoveWarp\).

\paragraph{Stress (Empty)}

\paragraph{Principal axes.}
With $\mathbf{U}_{\rm I}=\mathrm{diag}(d,-d)$, $\mathbf{V}_{\rm I}=\begin{bmatrix}0&d\\ d&0\end{bmatrix}$,
\[
u_I=\frac{y^2-z^2}{2d},\quad v_I=\frac{yz}{d}
\ \Longrightarrow\
\mathbf A
=\frac{y^2+z^2}{2}\,\mathbf{1}
+(\nu-1)\,\frac{y^2-z^2}{2}\,\mathrm{diag}(1,-1)
+(2\nu+1)\,yz
\begin{bmatrix}0&1\\[2pt]1&0\end{bmatrix},
\]
and $-\mathbf A\,\bm{\nu}$ reproduces the classical Love traction.

% \begin{equation}
% \left\{
% \begin{array}{ll}
% \Delta \bm{\chi}=0 & \text { in } \bar{\Section}, \\
% \nabla_{\bm{\nu}} \bm{\chi}=-\left[\frac{\nu}{2} \bm{R}+\left(1-\frac{\nu}{2}\right) \bm{Q}+\frac{2+\nu}{2} \bm{S}\right] \bm{\nu} & \text { on } \partial \bar{\Section}, \\
% \int_{\Section} \bm{\chi}\, \mathrm{d} A = \mathbf{0} & \text { (fix additive constants). }
% \end{array}
% \right.
% \end{equation}
% with $\bm{r}=(y,z)$ and $\bm{Q}=(\bm{r}\cdot\bm{r})\,\mathbf{1}-\bm{R}$ and 
% subject to the compatibility condition:
% \begin{equation}
% \int_{\Section} \bm{r}\, \mathrm{d} A=\mathbf{0}
% \end{equation}


This furnishes for \(\mathbf{e}_y\):
\begin{equation*}
\left\{
\begin{aligned}
\Delta \LoveWarp_y &= 0 \quad  &\text{in } \bar{\Section}, \\
\nabla_{\bm{\nu}}\LoveWarp_y &=-n_2\left[\frac{1}{2} \nu x_2^2+\left(1-\frac{\nu}{2}\right) x_3^2\right]-n_3(2+\nu) x_2 x_3
\quad&\text{on } \partial \bar{\Section}
\end{aligned}
\right.
\label{eq:love-flexure-2}
\end{equation*}

Stress is furnished by:
\begin{equation*}
\begin{aligned}
\tau_{xz} &= -\frac{V_z}{I}\frac{1}{2(1+\nu)}\left(
\chi_{,z} + \frac{1}{2}\nu z^2 + (1 - \tfrac{1}{2}\nu) y^2
\right) \\
\tau_{xy} &=
\end{aligned}
\end{equation*}
%

\subsubsection{Laplace-Neumann II (Sokolnikov)}
\cite{sokolnikoffMathematicalTheoryElasticity1946} introduced the harmonic re-gauge of Love's \eqref{eq:love-flexure-function}:
\begin{equation}
\left\{
\begin{aligned} 
& \Delta \varphi_1 = 0\\
& \frac{d \varphi_1}{d \nu}=\left[(1+\nu) x^2-\nu y^2\right] \cos (x, \bm{\nu}) 
+\left[(1+\nu) y^2-\nu x^2\right] \cos (y, \bm{\nu})
\end{aligned}
\right.
\end{equation}


\begin{equation}
\varphi_1=-\chi+1 / 3(1+1 / 2 \sigma)\left(x^3-3 x y^2\right)
\end{equation}
\begin{equation*}
\varphi_{\bm{p}}^{\text {Sok }}=-\chi_{\bm{p}}+\frac{1}{3}\left(1+\frac{\nu}{2}\right)(\nu) H_{\rm{p}}(\bm{r})+c_0+\mathbf{c} \cdot \bm{r}
\end{equation*}
Let the vector harmonic cubic be
$$
\HarmonicParticular(\bm{r})=\left(0, y^3-3 y z^2, z^3-3 z y^2\right)
\quad\text{with}\quad 
\Delta \HarmonicParticular=\mathbf{0}
$$
The directional scalar is $H_{\mathrm{p}}=\HarmonicParticular \cdot \bm{p}$.
Then the (directional) Sokolnikov potential is the harmonic re-gauge of Love:
$$
\varphi^{\mathrm{Sok}}=-\chi+\kappa_S(\nu) \mathbf{h}, \quad \kappa_S(\nu)=\frac{1}{3}\left(1+\frac{\nu}{2}\right) .
$$
with harmonic gauge \(h_{\rm p}\)
\begin{equation*}
h_{p}(\bm{r}):=(\bm{p} \cdot \bm{r})^3-3(\bm{p} \cdot \bm{r})\left((\mathbf{i}\times\bm{p}) \cdot \bm{r}\right)^2
\end{equation*}
\paragraph{Stress}
$$
\begin{aligned} & \tau_{x x}=\mu \alpha\left(\frac{\partial \varphi}{\partial x}-y\right)+\frac{V_z}{2(1+\nu) I_y}\left[\frac{\partial \varphi_1}{\partial x}-(1+\nu) x^2+\nu y^2\right] \\ & \tau_{x y}=\mu \alpha\left(\frac{\partial \varphi}{\partial y}+x\right)+\frac{V_z}{2(1+\nu) I_y} \frac{\partial \varphi_1}{\partial y} .\end{aligned}
$$

\subsubsection{Laplace-Neumann III (L0-N2, Serpieri-Rosati)}

Consider the torsionless vector warpage $\bm\psi:\bar\Omega\to\mathbb R^2$ defined by
\begin{equation}
\left\{
\begin{aligned}
\Delta \bm\psi &= \mathbf{0} && \text{in }\Omega,\\[2pt]
\partial_{\bm\nu}\bm\psi &= -\,\bm{A}^{\mathrm p} \bm\nu && \text{on }\partial\Omega,\\[2pt]
\displaystyle\int_{\Section}\bm\psi\mathrm{~d}A &=\mathbf{0}, \\ 
\int_{\Section}\bm{r}\mathrm{~d}A&=\mathbf{0},
\end{aligned}
\right.
\end{equation}
with
\begin{equation}
\bm{A}^{\mathrm p}
=\frac{1+\bar{\nu}}{4}\,\bm R+\frac{1-3\bar{\nu}}{8}\,s\,\bm 1
=\underbrace{\frac{1-\bar{\nu}}{4}}_{c_{\mathrm{iso}}}\,s\,\bm 1
+\underbrace{\frac{1+\bar{\nu}}{4}}_{c_{\mathrm{dev}}}\,\operatorname{dev}_{\Omega}\bm R
\;=\;c_{\mathrm{iso}}\,s\,\bm 1+c_{\mathrm{dev}}\,(u_I\,\bm U_I+v_I\,\bm V_I).
\end{equation}
Compatibility holds in centroidal axes because $\nabla\cdot\bm{A}^{\mathrm p}=\bm{r}$ and $\int_{\Section}\bm{r}\mathrm{~d}A=\mathbf{0}$.

\paragraph{Relation to Love's warpage}

Let $\bm\chi$ solve Love’s torsionless Laplace–Neumann problem \eqref{eq:love-flexure-function}. 
% \[
% \Delta\bm\chi=\mathbf{0},\qquad
% \partial_{\bm\nu}\bm\chi=-\,\bm{A}_{\mathrm L}\bm\nu,\qquad
% \int_{\Section}\bm\chi\mathrm{~d}A=\mathbf{0},
% \]
% with the compact boundary tensor
% \[
% \bm{A}_{\mathrm L}(\bm{r})
% =\tfrac12\,s\,\bm 1+(\nu-1)\,u_I\,\bm U_I+(2\nu+1)\,v_I\,\bm V_I.
% \]
Then
\[
\bm\psi=\bm\chi+\bm g_{L\to p},\qquad
\Delta \bm g_{L\to p}=\mathbf{0},\qquad
\partial_{\bm\nu}\bm g_{L\to p}=-(\bm{A}^{\mathrm p}-\bm{A}_{\mathrm L})\,\bm\nu,\qquad
\int_{\Section}\bm g_{L\to p}\mathrm{~d}A=\mathbf{0},
\]
where
\[
\bm{A}^{\mathrm p}-\bm{A}_{\mathrm L}
=\Big(\tfrac{1-\bar{\nu}}{4}-\tfrac12\Big)\,s\,\bm 1
+\Big(\tfrac{1+\bar{\nu}}{4}-(\nu-1)\Big)\,u_I\,\bm U_I
+\Big(\tfrac{1+\bar{\nu}}{4}-(2\nu+1)\Big)\,v_I\,\bm V_I.
\]

\paragraph{Relation to Poisson–Neumann warpage}

Let the bidirectional Poisson–Neumann warpage $\bm{\phi}$ be defined by
\[
\Delta\bm\phi=-2\,\bm{r},\qquad
\partial_{\bm\nu}\bm\phi=-\tfrac{\nu}{2}\,(\bm Q-\bm R)\,\bm\nu,\qquad
\int_{\Section}\bm\phi\mathrm{~d}A=\mathbf{0}.
\]
Recall the Poisson cubic vector field
\[
\PoissonParticular(\bm{r}) := v_I\,\bm V_I\,\bm{r},\qquad \Delta \PoissonParticular = 2\,\bm{r}.
\]
Then the classical componentwise identity generalizes to the vector form
\[
\bm{\phi} = -\,\bm{\chi} - \PoissonParticular \;+\; \bm c_0 + \bm C_1\,\bm{r},
\]
with an arbitrary affine harmonic gauge $(\bm c_0\in\mathbb R^2,\ \bm C_1\in\mathbb R^{2\times 2}$ skew). Consequently,
\[
\bm\psi = -\,\bm\phi + \PoissonParticular + \bm g_{\phi\to p},\qquad
\Delta \bm g_{\phi\to p}=\mathbf{0},\qquad
\partial_{\bm\nu}\bm g_{\phi\to p}=-\,\bm{A}^{\mathrm p}\bm\nu - \partial_{\bm\nu}(-\bm{\phi}+\PoissonParticular),\qquad
\int_{\Section}\bm g_{\phi\to p}\mathrm{~d}A=\mathbf{0}.
\]

\paragraph{Summary}

The problem with $\bm{A}^{\mathrm p}=c_{\mathrm{iso}}\,s\,\bm 1+c_{\mathrm{dev}}\,(u_I\,\bm U_I+v_I\,\bm V_I)$ is a Laplace–Neumann warpage with the same structure as Love’s, but assigning a single coefficient to the entire deviatoric content. It maps to Love’s harmonic warpage by a harmonic correction with quadratic Neumann flux equal to $\bm{A}^{\mathrm p}-\bm{A}_{\mathrm L}$, and to the Poisson–Neumann warpage by adding back the Poisson cubic $\bm f$ and a harmonic correction on the boundary.


\subsubsection{Poisson-Neumann I (L1-N2)}
The Poisson-Neumann problem for shear along \(\bm{p}\). 
\begin{equation}
\left\{
\begin{aligned}
\Delta (\bm{\phi}\cdot\bm{p}) &= -2 \;\bm{r}\cdot\bm{p} && \text{in } \bar{\Section},\\
\nabla_{\bm{\nu}}\,(\bm{\phi}\cdot\bm{p}) &=\nu (\operatorname{dev} \bm{r}\otimes\bm{r})\;\bm{p} \cdot\bm{\nu} && \text{on } \partial \bar{\Section}, \\
%=-\frac{\nu}{2} \left(\bm{Q}- \bm{R}\right) \bm{p}\cdot \bm{\nu}\\
\int_{\Section}\bm{r}\mathrm{~d}A &= \mathbf{0} && \text{integrability},\\
\int_{\Section}\bm{\phi}\mathrm{~d}A & = \mathbf{0} && \text{uniqueness}.
\end{aligned}
\right.
\end{equation}
This problem has been considered by \cite{romano2012torsion} with the symbol \(\phi^{\rm SH}\), \cite{pilkey2003analysis} with the symbols \(\Phi,\Psi\), 
\cite{lacarbonara2007solution} with the symbol \(\omega_i\),
and \cite{hjelmstad1983seismic} with the symbol \(\Psi\).

\paragraph{Stress}
\begin{equation*}
\bm{\tau} = G\left((\nabla \bm{\phi}^{\rm p})^T + \frac{\nu}{2}(\bm{Q} - \bm{R})\right)\bm{p}
= G\left((\nabla \bm{\phi}^{\rm p})^T - \nu \operatorname{dev} \bm{r}\otimes\bm{r}\right)\bm{p} 
\end{equation*}

Pilkey:
$$
\begin{aligned}
\tau_{x y} & =\frac{B}{2(1+\nu)}\left[\frac{\partial \Phi}{\partial y}+\nu\left(I_{y z} \frac{y^2-z^2}{2}-I_z y z\right)\right] \\
\tau_{x z} & =\frac{B}{2(1+\nu)}\left[\frac{\partial \Phi}{\partial z}+\nu\left(I_{y z} y z+I_z \frac{y^2-z^2}{2}\right)\right]
\end{aligned}
$$
$$
\begin{aligned}
&\begin{aligned}
\tau_{x y} & =\frac{V_y}{\Delta}\left(\frac{\partial \Psi}{\partial y}-d_y\right) \\
\tau_{x z} & =\frac{V_y}{\Delta}\left(\frac{\partial \Psi}{\partial z}-d_z\right)
\end{aligned}\\
\end{aligned}
$$
\begin{equation*}
\begin{aligned}
\tau_{12}&=\frac{V}{I_2}\frac{1}{2(1+\nu)}\left[\frac{\partial \HjelmstadPsi}{\partial x_2}-\frac{\nu}{2}\left(x_2^2-x_3^2\right)\right] \\
\tau_{13}&=\frac{V}{I_2}\frac{1}{2(1+\nu)}\left[\frac{\partial \HjelmstadPsi}{\partial x_3}-\nu x_2 x_3\right]
\end{aligned}
\end{equation*}

\paragraph{Relations}
\begin{equation}
\bm{\phi}=-(\bm{\chi} + \PoissonParticular)
\end{equation}
\iffalse
Any directional scalar used in the literature is just a dot with $\bm{p}$:

$$
\phi_{\rm p}=\bm{\phi} \cdot \bm{p}, 
\qquad \chi_{\rm p}=\boldsymbol{\chi} \cdot \bm{p}, 
\qquad f_{\rm p}=\PoissonParticular \cdot \bm{p},
$$
hence $\phi_{\mathrm{p}}=-\left(\chi_{\mathrm{p}}+f_{\mathrm{p}}\right)$ and everything is linear in \(\bm{p}\) .
$$
\Phi_{\bm{p}}=-\left(\chi_{\bm{p}}+f_{\bm{p}}\right)+c_0+\mathbf{c} \cdot \bm{r}+\left[\bm{r}^{\times}\right] \mathbf{d}+\left[\bm{r}^{\times}\right]^2 \mathbf{e} .
$$


On a simply connected section in centroidal axes $\int_{\Section} \bm{r} d A=\mathbf{0}$, the canonical choice that enforces
$\int_{\Section} \Phi_{\mathrm{p}}\, \mathrm{d} A=0$ and $\int_{\Section} \bm{r}\, \Phi_{\mathrm{p}}\, \mathrm{d} A=0$ is
$$
c_0=\frac{1}{A} \int_{\Section}\left(\chi_{\mathrm{p}}+f_{\mathrm{p}}\right) \, \mathrm{d} A, 
\quad 
\mathbf{c}=\mathbf{J}^{-1} \int_{\Section} \bm{r}\left(\chi_{\mathrm{p}}+f_{\mathrm{p}}\right) \,\mathrm{d} A, 
\quad \mathbf{d}=0, 
\quad \mathbf{e}=\mathbf{0} .
$$
To remove only the first moment along a particular direction \(\bm{p}\), use
$$
\mathbf{c}=\alpha \bm{p}, 
\quad \alpha=\frac{\int_{\Section}(\bm{p} \cdot \bm{r})\left(\chi_{\mathrm{p}}+f_{\mathrm{p}}\right) \,\mathrm{d} A}{\int_{\Section}(\bm{p} \cdot \bm{r})^2 \,\mathrm{d} A},
$$
i.e. divide by \(p\cdot \mathbf{J} p\).
\fi 

\paragraph{Principal Axes.}
Along a principal direction \(\mathbf{e}_y\), the form considered by \cite{hjelmstad1983seismic} is recovered:
\begin{equation*}
\left\{
\begin{aligned}
& \Delta \HjelmstadPsi =-2 x_2 \quad &\text{in } \Omega \\
& \nabla_{\bm{\nu}} \HjelmstadPsi=\frac{\nu}{2}\left[\left(x_2^2-x_3^2\right) n_2+2 x_2 x_3 n_3\right]
\end{aligned}
\right.
\end{equation*}
%
\begin{equation*}
\HjelmstadPsi = - (\LoveWarp + x_2 x_3^2) + \text{const.}
\end{equation*}
%
\begin{Quote}

\cite{lacarbonara2007solution} for \(j=1,2,3,\):
\begin{equation*}
\left\{
\begin{aligned}
\Delta \omega_j&=-2\left(\frac{x_1}{J_2} \delta_{1 j}+\frac{x_2}{J_1} \delta_{2 j}\right) &\text { in } \Omega ; \\
\nabla_{\bm{\nu}}\,\omega_j&=g_j &\text { on } \partial \Omega,  \\
%
g_1&=\frac{\nu}{J_2}\left[\frac{1}{2}\left(x_1^2-x_2^2\right) n_1+x_1 x_2 n_2\right], \\
g_2&=-\frac{\nu}{J_1}\left[\frac{1}{2}\left(x_1^2-x_2^2\right) n_2-x_1 x_2 n_1\right], \\ 
g_3&=x_2 n_1-x_1 n_2,
\end{aligned}
\right.
\end{equation*}
Stresses are:
\begin{equation*}
\bm{\tau}=\frac{G}{E}\left(Q_1 \nabla \omega_1+Q_2 \nabla \omega_2\right)
+G\left[k\left(\nabla \omega_3+\mathbf{a} \times \mathbf{x}\right)+\nu \mathbf{v}\right] .
\end{equation*}
\begin{equation*}
\begin{aligned}
& \mathbf{v}_{\varepsilon}(\bm{x})=-\varepsilon\left(x_1 \mathbf{a}_1+x_2 \mathbf{a}_2\right) \\
& \mathbf{v}_\kappa(\bm{x})=\left[-k_1 x_1 x_2+\frac{1}{2} k_2\left(x_1^2-x_2^2\right)\right] \mathbf{a}_1+\left[k_2 x_1 x_2+\frac{1}{2} k_1\left(x_1^2-x_2^2\right)\right] \mathbf{a}_2, \\
& \mathbf{v}_\mu(\bm{x})=\left[-\mu_1 x_1 x_2+\frac{1}{2} \mu_2\left(x_1^2-x_2^2\right)\right] \mathbf{a}_1+\left[\mu_2 x_1 x_2+\frac{1}{2} \mu_1\left(x_1^2-x_2^2\right)\right] \mathbf{a}_2
\end{aligned}
\end{equation*}
\end{Quote}

\begin{Quote}
% \paragraph{Pilkey}
\cite{pilkey2003analysis} uses this form of the problem with the notation 
\begin{equation*}
    \begin{aligned}
        \Phi &= \bm{\phi}\cdot \mathbf{I}\,\mathbf{i}_z \\
        \Psi &= \bm{\phi}\cdot \mathbf{I}\,\mathbf{i}_y\\
        \mathbf{h} &= \bm{\Gamma}\,\mathbf{I}\,\mathbf{i}_z \\
        \mathbf{d} &= \bm{\Gamma} \, \mathbf{I}\, \mathbf{i}_y
    \end{aligned}
\end{equation*}
Thus:
\begin{equation}
\begin{aligned}
\Delta \Phi&=2\left(I_{y z} y-I_z z\right)  \\
   \bm{\nu} \cdot \nabla \Phi&=\bm{\nu}\cdot \mathbf{h}
\end{aligned}
\qquad\text{and}\qquad
\begin{aligned}
\Delta \Psi&=2\left(I_{y z} z-I_y y\right) \\
\bm{\nu} \cdot \nabla \Psi&=\bm{\nu} \cdot \mathbf{d} \\
\end{aligned}
\end{equation}
with
\begin{equation}
\begin{aligned}
\mathbf{h}=\nu\left(I_z y z-I_{y z} \frac{y^2-z^2}{2}\right) \mathbf{j}-\nu\left(I_{y z} y z+I_z \frac{y^2-z^2}{2}\right) \mathbf{k} \\
\mathbf{d}=\nu\left(I_y \frac{y^2-z^2}{2}-I_{y z} y z\right) \mathbf{j}+\nu\left(I_y y z+I_{y z} \frac{y^2-z^2}{2}\right) \mathbf{k}
\end{aligned}
\end{equation}

\begin{equation}
\begin{aligned}
\Delta \bm{\psi}=2 \nu \mathbf{I}\bm{r} \quad \text { in } \Omega \\
\partial_\nu \psi=\nu \mathbf{I}\left(\left(\operatorname{dev} \bm{R}\cdot \mathbf{C}_I\right) \frac{1}{2} \mathbf{1}+\left(\operatorname{dev} \bm{R}\cdot \mathbf{B}_I\right) \mathbf{J}\right) \bm{\nu} &
\end{aligned}
\end{equation}
\end{Quote}


\subsubsection{Poisson-Neumann II (L1-N2,  Irrotational, Gruttmann)}
\begin{equation}
\left\{
\begin{aligned}
\Delta \bar{\GruttmannPsi}&=-\bm{r}\cdot{\bm{a}} & \text { in } \Omega \\
\nabla_{\bm{\nu}}\bar{\GruttmannPsi}&=n_y g_1(z)-n_z g_2(y)
\end{aligned}
\right.
\end{equation}
with
\begin{equation*}
\begin{aligned}
f_1(y, z)&=a_1 \bar{y}+a_2 \bar{z} = \bm{r}\cdot\bm{a} \\
g_1(z)&=-\frac{1}{2} \bar{\nu}\, a_1\left(z-z_0\right)^2 \\
g_2(y)&=\frac{1}{2} \bar{\nu}\, a_2\left(y-y_0\right)^2
\end{aligned}
\end{equation*}

\begin{equation*}
\LoveWarp\left(x_2, x_3\right)=\frac{1}{G} \bar{\Psi}\left(x_2, x_3\right)+\bm{c} \cdot (\bm{r} + \mathbf{i}),
\end{equation*}
\begin{equation}
\begin{aligned}
\tau_{x y} &=\bar{\Psi},_y-g_1(z) \\
\tau_{x z} &=\bar{\Psi},_z+g_2(y)
\end{aligned}
\end{equation}
\begin{equation}
\tau_{x y}=\bar{\mathbf{\Psi}}_{1, y}+\mathbf{\Psi}_{2, z}, 
\quad \tau_{x z}=\bar{\mathbf{\Psi}}_{1, z}-\mathbf{\Psi}_{2, y},
\end{equation}

Poisson correction (torsionless): \[
\psi_P=\frac{\bar{\GruttmannPsi}_2}{G}
=c_x \varphi_1+c_y \varphi_2
\]
$c_x, c_y$ depend on the chosen shear/bending normalization; pick them so that $\nabla_{\bm{\nu}} \psi_P$ matches the \(\nu\)-dependent quadratic part of the boundary condition.


\subsubsection{Poisson-Neumann III (L1-N2, Simo)}
Define the scaling:
\begin{equation*}
\begin{aligned}
&\mathbf{C} \triangleq -\kappa \frac{G A}{E} \mathbf{I}^{-1},
\end{aligned}
\end{equation*}
Find functions \(\bm{\varphi}\) and Lagrange multipliers \(\bm{\Lambda}\) that satisfies:
\begin{equation}
\left\{
\begin{aligned}
\Delta \bm{\varphi} & =2 \mathbf{C} \bm{r} & & \text { in } \bar{\Section}, \\
\nabla_{\bm{\nu}} \bm{\varphi} & =-\frac{\nu}{2} \mathbf{C}\bigl(\bm{R} - \bm{Q}\bigr) \bm{\nu} + \bm{\Lambda} \bm{\nu} & & \text { on } \partial \bar{\Section}, \\
\int \bm{r}\mathrm{~d}A &= \mathbf{0} & & \text{integrability} \\
\int \bm{\varphi}\mathrm{~d}A &= \mathbf{0} & & \text{uniqueness}, \\
\int_\Omega \bm{r}\otimes \boldsymbol{\varphi}\mathrm{~d}A&=\mathbf{0}  & & \text{first moment orthogonality constraints.}
\end{aligned}
\right.
\end{equation}
%
%
Here $\bm{\Lambda}$ is a ${2} \times {2}$ matrix of Lagrange multipliers that would arise from adding a harmonic affine field $\bm{c}+\mathbf{B} \bm{r}$ to $\varphi$ (since $\Delta\,(\bm{c}+\mathbf{B} \bm{r})=\mathbf{0}$ and $\partial_\nu(\bm{c}+\mathbf{B} \bm{r})=\mathbf{B} \bm{\nu}$). 
The constants $\bm{c}$ enforce the zero-mean, the matrix $\mathbf{B}$ enforces the ``moment" constraints.


Consider the projection functional that removes mean and all first moments:
\begin{equation}
\Pi[\bm{u}]:=\bm{u}-\frac{1}{A} \int \bm{u} \, \mathrm{d} A - \left(\mathbf{J}^{-1} \int \bm{r} \otimes \bm{u}\, d A\right) \bm{r}
\end{equation}
\begin{equation}
\bm{\varphi}=\mathbf{D} \Pi[\bm{\psi}], 
\quad\text{with}\quad
\mathbf{D}= \frac{GA}{E} \mathbf{I}^{-1}\bm{\Xi}= \frac{A}{2(1+\nu)} \mathbf{I}^{-1}\bm{\Xi}
\end{equation}

\begin{equation}
\varphi_{\mathrm{e}}^{\text {Simo }}=\varphi^{\text {Simo }} \cdot \bm{p}= \kappa \frac{A}{2(1+\nu)} \left(\mathbf{I}^{-1} \bm{p}\right) \cdot \mathbf{\Pi}[\mathbf{\Phi}],
\end{equation}
$$
\int_{\Section}\left(\bm{\SimoWarp} \cdot \bm{\gamma}\right) \,\mathrm{d} A
=\frac{\kappa G A}{E} \int_{\Section} \Pi[\bm{\HjelmstadPsi}] \cdot\left(\mathbf{I}^{-1} \bm{\gamma}\right) \,\mathrm{d} A 
\stackrel{*}{=} \int_{\Section} \bm{\HjelmstadPsi} \cdot \bm{p} \,\mathrm{d} A,
$$
where $\bm{p}=\mathbf{J}^{-1} \bm{q}/E$ and $*$ uses:
\begin{itemize}
\item $\int_{\Section} \Pi[\bm{\HjelmstadPsi}]\, \mathrm{d} A=\mathbf{0}$,
\item $\int_{\Section} \bm{r} \,\mathrm{d} A=\mathbf{0}$, so $\int \bm{a} \cdot \Pi[\Phi]=\int \bm{a} \cdot \Phi$ for any constant vector $\bm{a}$ (which $\mathbf{I}^{-1} \bm{\gamma}$ is under Timoshenko).
\end{itemize}
Thus contraction of \(\bm{\SimoWarp}\) with $\bm{\gamma}$ is energetically equivalent to contracting $\bm{\HjelmstadPsi}$ with $\bm{p}$. 
Pointwise equality is not expected in general rotated frames.

\paragraph{Stress}
$$
\begin{aligned}
& \sigma_{12}=-\left(\frac{G \Omega}{E I_2}\right) \frac{V}{\Omega}\left[\phi, 2+\frac{\nu}{2}\left(x_2\right)^2+\left(1-\frac{\nu}{2}\right)\left(x_3\right)^2\right] \\
& \sigma_{13}=-\left(\frac{G \Omega}{E I_2}\right) \frac{V}{\Omega}\left[\phi, 3+(2+\nu) x_2 x_3\right]
\end{aligned}
$$
\paragraph{Principal Axes.}
In principal axes this collapses to the componentwise formula of \cite{simo1982consistent}:
$$
\varphi_y=\mathbf{C}_{y y}\left(\Phi_y-\left\langle\Phi_y\right\rangle-\frac{y}{J_{y y}} \int_{\Section} y \Phi_y d A\right), \quad \varphi_z=\mathbf{C}_{z z}\left(\Phi_z-\left\langle\Phi_z\right\rangle-\frac{z}{J_{z z}} \int_{\Section} z \Phi_z d A\right)
$$
\begin{equation*}
\varphi_y(y, z)=-\kappa \frac{G A}{E I_z}\left[f-\underbrace{\frac{1}{A} \int_{\Section} f \,\mathrm{d} \Omega}_{\text {remove mean }}-\underbrace{\frac{y}{I_z} \int_{\Section} y f \,\mathrm{d} \Omega}_{\text {remove } y \text {-moment }}\right]
\end{equation*}

\begin{equation}
\begin{aligned}
\SimoWarp_y
&=-\kappa\frac{G A}{E I_z}\left[\bigl(\LoveWarp_y + y z^2\bigr)-\frac{1}{A} \int_{\Section}\left(\chi_y+y z^2\right)\, \mathrm{d} \Omega-\frac{y}{I_z} \int_{\Section} y\left[\chi+y z^2\right]\, \mathrm{d} \Omega\right] \\
&= \kappa\frac{A}{2(1+\nu) I_z}\left[\HjelmstadPsi -\frac{1}{A} \int_{\Section} \HjelmstadPsi\, \mathrm{d} \Omega
-\frac{y}{I_z} \int_{\Section} y \HjelmstadPsi \, \mathrm{d} \Omega\right] \\
\end{aligned}
\end{equation}



\begin{Quote}
% Enforce $\int_{\Section}(\mathbf{e} \cdot \bm{r}) \varphi_{\mathrm{e}} \,\mathrm{d} A=0$ by choosing $\Lambda$ full
% \begin{equation*}
% \lambda_2=-\frac{\kappa GA}{EI_2} \frac{1}{I_2} \int_{\Section} x_2 \HjelmstadPsi\, \mathrm{d} A
% \end{equation*}

Let $\bm{r}=(y, z)$ in centroidal coordinates $\int_{\Section} \bm{r} \,\mathrm{d} A=\mathbf{0}$.
Solve the Simo BVP once with $\bm{\Lambda}=\mathbf{0}$ to get a provisional field $\bm{\varphi}^0$ (then remove its mean so $\int_{\Section} \bm{\varphi}^0 \,\mathrm{d} A=\mathbf{0}$ ). 
The boundary multiplier just injects a linear gauge:
$$
\bm{\varphi}=\bm{\varphi}^0+\bm{\Lambda}\bm{r}, 
\quad 
\nabla_{\bm{\nu}} \bm{\varphi}=\nabla_{\bm{\nu}} \bm{\varphi}^0+\bm{\Lambda} \bm{\nu},
$$
so the $\bm{\Lambda}\bm{\nu}$ term in the BC is exactly the linear correction.
Moments tensor and mixed moment
% Recall the $2 \times 2$ second-moment tensor
% %
% $$
% \mathbf{J}=\int_{\Section} \bm{r} \otimes \bm{r}\, \mathrm{d} A
% $$
%
The mixed moment (a $2 \times 2$ matrix) is:

$$
\mathbf{M}^0=\int_{\Section} \bm{r} \otimes \bm{\varphi}^0 \,\mathrm{d} A=\left[\begin{array}{lll}
\int y \varphi_2^0 & \int y \varphi_3^0 \\
\int z \varphi_2^0 & \int z \varphi_3^0
\end{array}\right]
$$
Enforce $\int_{\Section} \bm{r} \otimes \bm{\varphi} \,\mathrm{d} A=\mathbf{0}$.
Since $\int_{\Section} \bm{r} \otimes(\bm{\Lambda}\bm{r}) \,\mathrm{d} A=\mathbf{J} \mathbf{\Lambda}^{\rm t}$,

$$
\mathbf{\Lambda}^{\top}=-\mathbf{J}^{-1} \mathbf{M}^0 
\quad \Longleftrightarrow \quad 
\bm{\Lambda}=-\mathbf{M}^{0^{\top}} \mathbf{J}^{-1}
$$
\end{Quote}





\subsection{Dirichlet Problems}
\subsubsection{Poisson-Dirichlet (L1-D0, Potential Stress)}
\cite{lacarbonara2007solution} (who used \(\psi\))
\begin{equation}
\left\{
\begin{aligned}
\nabla^2 \hat{\Psi}_j&=\bar{\nu}\left(\frac{x_2}{J_2} \delta_{1 j}-\frac{x_1}{J_1} \delta_{2 j}\right)-2 \delta_{3 j} \text { in } \Omega ; \\
\hat{\Psi}_j&=\bar{g}_j(\bm{r}) \text { on } \partial \Omega\\
\text { with }&\\
\bar{g}_1&=-\frac{1}{2 J_2} \int_0^s x_1^2 \mathrm{~d} x_2, \\
\bar{g}_2&=\frac{1}{2 J_1} \int_0^s x_2^2 \mathrm{~d} x_1, \\
\bar{g}_3&=0
\end{aligned}
\right.
\end{equation}

\subsubsection{Poisson-Dirichlet (L1-D3, Solenoidal Stress, Gruttmann)}
\cite{gruttmann1999shear}:
\begin{equation}
\left\{
\begin{aligned}
\Delta \Psi=-f_2(y, z) \quad &\text{in } \Omega \\ \Psi(s)=h_1(s) \quad &\text{on } \partial \Omega
\end{aligned}
\right.
\end{equation}
\begin{equation*}
f_2(y, z)=\bar{\nu}\left[a_2\left(y-y_0\right)-a_1\left(z-z_0\right)\right]
\end{equation*}
\begin{equation*}
h_1(s)=\Psi_0+\frac{1}{2}\left(a_1 \bar{y}^2 \bar{z}-a_2 \bar{z}^2 \bar{y}\right)
\end{equation*}
\begin{equation*}
\begin{aligned}
\tau_{x y} & = \Psi,_z-\frac{1}{2} a_1 \bar{y}^2 \\
\tau_{x z} & =-\Psi,_y-\frac{1}{2} a_2 \bar{z}^2
\end{aligned}
\end{equation*}
\begin{equation}
\Psi=-J_2 a_2 \psi_1-J_1 a_1 \psi_2+\Psi_0
\end{equation}

\subsubsection{Poisson-Dirichlet (L2-D1, Total Potential, Renton)}

$\phi$ is a stress function which satisfies the equation
\begin{equation}
\left\{
\begin{aligned}
\Delta \phi=\bar{\nu} \frac{V}{I}y-\frac{\mathrm{d} f}{\mathrm{~d} y} \quad& \text{in } \Omega\\
\phi = \text{const.} \quad&\text{on }\partial\Omega
\end{aligned}
\right.
\end{equation}
on the cross-section, $f(y)$ is a function such that
$$
\frac{\partial \phi}{\partial s}=\left[\frac{V}{2 I} x^2-f(y)\right] \frac{\mathrm{d} y}{\mathrm{~d} s}=0
$$
on the boundary of the cross-section, where $s$ is measured around the boundary, $x$ and $y$ are the principal axes of the cross-section through its centroid and $I$ is the second moment of area of the cross-section about the $y$ axis. Without loss of generality, $\phi$ can be taken as identically zero on the boundary.

$$
\tau_{y z}=-\frac{\partial \phi}{\partial x}
$$


\subsubsection{Poisson-Neumann (Displacement)}
\begin{equation}
\begin{aligned}
\Delta \psi&=-\frac{1}{G} f_1, \\
\partial_n \psi&=-\left(c_1 n_2+c_2 n_3\right)
\end{aligned}
\end{equation}



Let $\bar{\psi}$ be the gauge-fixed flexural warping (i.e., $\bar{\psi}=\psi+c_2 x_2+c_3 x_3$ ). Then for shear along $x_2$,

$$
\Delta \bar{\psi}=-\frac{1}{G} a_1 x_2 \quad \text { in } \Omega, \quad \partial_n \bar{\psi}=0 \text { on } \partial \Omega
$$

and $\bar{\Psi}_1=G \bar{\psi}$. The Poisson piece $\psi_P$ is harmonic with the Neumann data shown above, $\psi_P=\bar{\Psi}_2 / G$. The total torsionless warping is $\psi_{\text {tot }}=\bar{\psi}+\psi_P$ (plus an arbitrary affine gauge).


\subsection{Shear Corrections}

\cite{puchegger2003experimental} found that the $\kappa$-factor decreases with increasing depth-to-width ratio, and that dependence on the aspect ratio increases with increasing Poisson’s ratio .

The shear coefficient of the deep cross-sections is notably higher than the one associated with the shallow cross-sections, and accordingly, the shear coefficient is anticipated to designate an increasing function in terms of the aspect ratio.

The shear flexibility of a rectangular section with a constant cross-sectional area increases in proportion to $(b / a)^2$ as the section becomes a very thin, flat strip \citep{renton1991generalized}.
However, since the bending flexibility of the section will increase in proportion to $b / a$, the result is not so surprising.

\subsubsection{Setup}
Write the exact (torsionless) stress as

$$
\boldsymbol{\tau}(\bm{r})=\mathbf{T}(\bm{r}) \bm{p}, \quad \bm{r}=(y, z),
$$
where $\mathbf{T}(\bm{r}) \in \mathbb{R}^{2 \times 2}$ is the section kernel you get from whatever Saint-Venant BVP you adopt (Hjelmstad/Romano/Pilkey, etc.). It depends on geometry (and $\nu$ ), but not on $G$. Local shear strain is $\bm{\varepsilon}_{\gamma}(\bm{r})=\bm{\tau} / G=\frac{1}{G} \mathbf{T} \bm{~p}$.

\subsubsection{Energetic (Renton)}
\cite{gruttmann2001shear} adopted the following definition, which had been introduced by \cite{bach1924elastizitaet} and Stojek (1964):
\begin{equation*}
\frac{V_x^2}{A_x}+\frac{V_z^2}{A_z} = \SectionIntegral{\sigma_{y x}^2+\sigma_{y z}^2}
\end{equation*}
where the shear-corrected areas, $A_x$ and $A_z$, are computed by letting $V_z$ and $V_x$ be equal to zero, respectively. 
More generally, \citep{renton1991generalized,serpieri2014equivalence}
$$
\mathbb{I}\bm{n} \cdot \bm{D}_{\mathrm{v}}  \mathbb{I}\bm{n} 
=
\frac{1}{G} \int_0^l \int_{\Section} \bm{\tau}_{\mathrm{v}}^{(D S V)} \cdot \bm{\tau}_{\mathrm{v}}^{(D S V)} d A d z
$$
furnishing
$$
\bm{D}_{\mathrm{v}}=\frac{L}{G} \mathbf{J}_G^{-1} \bm{\Phi}^{(E)} \mathbf{J}_G^{-1},
\quad \text { where } \quad
\bm{\Phi}=\SectionIntegral{\left(\nabla_k \bm{\psi} \otimes dX^k +\bm{A}^{\mathrm{p}}\right)\left(dX^k \otimes \nabla_k \bm{\psi}+\bm{A}^{\mathrm{p}}\right)}
$$

Renton matches the total shear energy per unit length for any $\bm{p}$ :

$$
\begin{gathered}
U_{\mathrm{sv}}(\mathbf{p})=\frac{1}{2 G} \int_{\Section}\|\boldsymbol{\tau}\|^2 d A=\frac{1}{2 G} \bm{p}\cdot\left(\int_{\Section} \mathbf{T}^{\top} \mathbf{T} d A\right) \bm{p}, \\
U_{\text {beam }}(\mathbf{p})=\frac{1}{2} \mathbf{p}^{\top} \gamma_{\text {beam }}=\frac{1}{2 G A} \mathbf{p}^{\top} \mathbf{\Xi}^{-1} \mathbf{p} .
\end{gathered}
$$

Equating for all $\mathbf{p}$ yields

$$
\Xi_{\text {Renton }}^{-1}=A \int_{\Section} \mathbf{T}^{\top}(\bm{r}) \mathbf{T}(\bm{r}) d A, \quad \Xi_{\text {Renton }}=\left(A \int_{\Section} \mathbf{T}^{\top} \mathbf{T} d A\right)^{-1}
$$
\begin{equation}
\frac{1}{\kappa}=\frac{1}{G I^2} \int_A\left(\frac{I}{V}f-\frac{1}{2} x^2\right)^2-\phi\, \frac{I}{V}\left(\frac{I}{V}\frac{\mathrm{d} f}{\mathrm{~d} y}+\frac{\nu}{1+\nu}y\right) \mathrm{d} A .
\end{equation}

\subsubsection{Geometric (Cowper)}
\cite{cowper1966shear} considered the mean deflection of the cross-section \(\tilde{u}_z\), the mean angle of rotation of the cross-section around the neutral axis $\Theta$ and the total transverse shear force \(V\)acting on the cross-section,
$$
\begin{gathered}
\tilde{u}_z=\frac{1}{A} \SectionIntegral{ u_z } \\
\Theta = \frac{1}{I} \SectionIntegral{ x u_y } \\
V=\SectionIntegral{ \sigma_{x y} }
\end{gathered}
$$
where \(A\) is the cross-section area, and \(I\) is the moment of inertia of the cross-section. The shear-corrected area $A_{\mathrm{v}}$ is computed from
$$
\frac{\partial \tilde{u}_z}{\partial y}+\Theta = \frac{V}{A_{\mathrm{v}} G}
$$

\begin{equation}
\kappa_y=\frac{1}{1+\frac{1}{A} \int \frac{\partial \SimoWarp}{\partial y} \, \mathrm{d} \Omega}
=\frac{2\left(1+\nu\right) I_z}{\nu \frac{I_y-I_z}{2}-\frac{A}{I_z} \int_{\Section} y\left[\LoveWarp+y z^2\right] d \Omega}
\end{equation}

\subsubsection{Hutchinson}
\cite{hutchinson2001shear}
\begin{equation}
\kappa=-\dfrac{2(1+\nu)}{\dfrac{A}{I_z^2} C_4+\nu\left(1-\dfrac{I_y}{I_z}\right)}
\end{equation}
\begin{equation*}
C_4=-\int_A\left[\nu\left(f_1 y^2-f_1 z^2+2 f_2 y z\right)+2(1+\nu)\left(f_1^2+f_2^2\right)\right]\,\mathrm{d} A
\end{equation*}
\begin{equation}
\begin{gathered}
f_1=-\frac{1}{2(1+\nu)}\left(\frac{\partial \LoveWarp}{\partial y}+\frac{\nu y^2}{2}+\frac{(2-\nu) z^2}{2}\right) \\
f_2=-\frac{1}{2(1+\nu)}\left(\frac{\partial \LoveWarp}{\partial z}+(2+\nu) y z\right)
\end{gathered}
\end{equation}

\subsection{Examples}

